\documentclass[11pt,a4paper,sans]{moderncv}
\moderncvstyle{classic}
\moderncvcolor{blue}
\usepackage[utf8]{inputenc}
\usepackage[scale=0.8]{geometry}
\usepackage{helvet}

\name{Alexandro}{Disla}
\title{Spécialiste MEAL | Analyste de Données \& Développeur Backend}
\address{Rue Saint-Louis Jeanty \#14, Route de Frère}{}{}
\phone[mobile]{(+509) 4148-3700}
\email{alexandrodisla@hotmail.com}
\extrainfo{GitHub: \url{https://github.com/AD0791}}

\begin{document}
\makecvtitle

\section{Profil Professionnel}
\cvitem{}{Professionnel orienté résultats avec plus de 5 ans d'expérience en Suivi, Évaluation, Redevabilité et Apprentissage (MEAL) et systèmes de données. Expert dans la conception de cadres MEAL, l'automatisation de pipelines de données et la création d'outils de visualisation avancés pour soutenir la prise de décision fondée sur des preuves dans des contextes humanitaires et de développement. Expérience confirmée dans l'alignement des indicateurs de projet avec les cadres logiques et les exigences des bailleurs (Rotary International, USAID).}

\cvitem{Points forts}{
\begin{itemize}
    \item \textbf{Conception de systèmes MEAL :} Élaboration de plans MEAL, cadres logiques et tableaux de suivi des indicateurs (ITT) ;
    \item \textbf{Automatisation et intégration de données :} SQL avancé, ETL, Python et outils de collecte numérique (Commcare, mWater, Kobo) ;
    \item \textbf{Reporting et visualisation :} Expert Power BI, Looker Studio, RMarkdown et Quarto pour un reporting prêt pour les bailleurs ;
    \item \textbf{Qualité et analyse des données :} Mise en œuvre de protocoles DQA et réalisation d'analyses de tendances complexes ;
    \item \textbf{Polyvalence technique :} Développement d'API sécurisées (FastAPI, Docker) et infrastructure cloud (AWS, GCP) ;
    \item \textbf{Base analytique :} Formation en économie appliquée avec un focus sur la recherche quantitative.
\end{itemize}}

\section{Expériences Professionnelles}
\cventry{Févr 2026 -- Présent}{Ingénieur Logiciel}{Tekkod LLC}{Télétravail}{}{
\begin{itemize}
    \item Direction du projet << STS Mobile App Modernization >>, mise à niveau de React Native de 0.65.1 vers 0.74.x.
    \item Migration de l'architecture vers le moteur Hermes, Java 17 et Gradle 8.8, activant le mode Bridgeless.
    \item Migration complète vers Firebase v22 (API modulaire) et remplacement des notifications obsolètes par @notifee/react-native.
    \item Audit UI Yoga 3.0 pour assurer la cohérence du layout et la compatibilité avec React 18.
    \item Refactorisation des composants de base et correction des mises à jour d'état asynchrones dans RewardsModal.
    \item Mise à jour du dashboard Looker Studio pour ajouter des capacités de recherche et de filtrage.
\end{itemize}}

\cventry{Oct 2023 -- Jan 2026}{Consultant mWater \& Analyste de Données}{HANWASH}{Télétravail (Upwork)}{}{
\begin{itemize}
    \item \textbf{Conception de système MEAL :} Participation active à la conception et au raffinement du plan MEAL du projet, assurant l'alignement avec le cadre logique et les normes JMP.
    \item \textbf{Suivi des indicateurs :} Développement et maintenance du tableau de suivi des indicateurs (ITT) pour monitorer les progrès et ajuster les interventions.
    \item \textbf{Gestion des données :} Conception et déploiement d'enquêtes mWater complexes ; résolution des problèmes de collecte pour garantir l'intégrité des données.
    \item \textbf{Visualisation :} Création de tableaux de bord automatisés (Looker Studio, Power BI, Plotly) pour le reporting aux parties prenantes.
    \item \textbf{Renforcement des capacités :} Production de guides techniques et formation des équipes de terrain sur la collecte numérique et les protocoles qualité.
\end{itemize}}

\cventry{Mars 2021 -- Mai 2024}{Développeur Backend}{Tekkod LLC}{}{}{
\begin{itemize}
    \item Contrat à temps partiel. Développement de services web RESTful conformes aux bonnes pratiques d'architecture logicielle.
    \item Mise en place de conteneurs Docker et pipelines de déploiement pour un environnement stable.
    \item Sécurisation via JWT et chiffrement des mots de passe.
    \item Application du repository pattern pour séparer les couches contrôleur et base de données.
    \item Documentation complète des endpoints et gestion des versions via Git.
    \item Mise en place d'un data pipeline assurant la migration de la base MongoDB vers BigQuery (Google Cloud) via Apache Beam.
    \item Création et maintenance d'un dashboard analytique sur Looker Studio basé sur les données BigQuery.
\end{itemize}}

\cventry{Oct 2021 -- Août 2024}{Monitoring \& Evaluation Officer}{Impact Youth Project, Caris Foundation International}{}{}{
\begin{itemize}
    \item \textbf{Mise en place de systèmes :} Conception et implémentation de systèmes de collecte numérique (Commcare, HIVHaiti), optimisant les flux de données multi-sites.
    \item \textbf{Reporting automatisé :} Développement de rapports d'activité et de bailleurs automatisés (Python, SQL), réduisant les délais de production.
    \item \textbf{Assurance Qualité (DQA) :} Établissement et supervision des protocoles DQA, garantissant la fiabilité des données intégrées sous MySQL.
    \item \textbf{Déploiement de Dashboards :} Création et déploiement de tableaux de bord hebdomadaires via AWS (RMarkdown, Quarto) pour le suivi des KPIs.
    \item \textbf{Suivi de performance :} Maintenance du dashboard Power BI pour le programme EID, avec analyses de tendances bi-hebdomadaires.
\end{itemize}}

\cventry{Nov 2015 -- Mars 2021}{Consultant en Développement Logiciel (Projet UNOPS)}{CassionSoft}{}{}{
\begin{itemize}
    \item Création d'un service REST API pour l'analyse de formulaires, avec authentification OAuth2 / JWT.
    \item Implémentation de la stratégie d'autorisation et documentation technique complète.
\end{itemize}}

\cventry{2014 -- 2015}{Économiste -- Service de Planification Économique}{Ministère de la Planification et de la Coopération Externe}{}{}{
\begin{itemize}
    \item Suivi des indicateurs macroéconomiques et participation à la rédaction de rapports de politiques publiques.
    \item Collaboration avec l'équipe de modélisation pour les simulations d'impact socio-économique.
\end{itemize}}

\section{Formation}
\cvitem{2010--2014}{Diplôme d'Études Sup\'{e}rieures en Économie Appliquée, C.T.P.E.A}
\cvitem{2009--2010}{Baccalauréat (Philo, mention Bien), Institution Saint-Louis de Gonzague}
\cvitem{Autoformation continue}{FlatIronSchool, CodeWithMosh, Udemy}

\section{Compétences Techniques}
\cvitem{Langages \& Frameworks}{Python, FastAPI, SQLModel, SQLAlchemy, Pydantic, Pandas, Numpy, Plotly, Docker, Git, React, React Native}
\cvitem{Bases de données}{MySQL, PostgreSQL, SQL Server, MongoDB (Beanie, Pymongo), BigQuery}
\cvitem{Outils \& Automatisation}{Commcare, Streamlit, Power BI, Looker Studio, RMarkdown, Quarto, Apache Beam, AWS, Google Cloud, VS Code}
\cvitem{Méthodologies}{CI/CD, API Documentation, Unit Testing, Monitoring \& Métriques, Intégration de données (ETL)}
\cvitem{Langues}{Français (courant), Anglais (courant), Créole (natif)}

\section{Points Distinctifs}
\cvitem{}{Capacité à comprendre rapidement les processus métiers et à les traduire en solutions techniques.}
\cvitem{}{Forte orientation vers la qualité du code, la sécurité et la performance.}
\cvitem{}{Esprit d'équipe, rigueur et curiosité technologique continue (IA, automatisation, Power Platform).}

\end{document}
